% Options for packages loaded elsewhere
\PassOptionsToPackage{unicode}{hyperref}
\PassOptionsToPackage{hyphens}{url}
\PassOptionsToPackage{dvipsnames,svgnames,x11names}{xcolor}
%
\documentclass[
  letterpaper,
  DIV=11,
  numbers=noendperiod]{scrartcl}

\usepackage{amsmath,amssymb}
\usepackage{iftex}
\ifPDFTeX
  \usepackage[T1]{fontenc}
  \usepackage[utf8]{inputenc}
  \usepackage{textcomp} % provide euro and other symbols
\else % if luatex or xetex
  \usepackage{unicode-math}
  \defaultfontfeatures{Scale=MatchLowercase}
  \defaultfontfeatures[\rmfamily]{Ligatures=TeX,Scale=1}
\fi
\usepackage{lmodern}
\ifPDFTeX\else  
    % xetex/luatex font selection
    \setmonofont[Scale=0.75]{Latin Modern Mono}
\fi
% Use upquote if available, for straight quotes in verbatim environments
\IfFileExists{upquote.sty}{\usepackage{upquote}}{}
\IfFileExists{microtype.sty}{% use microtype if available
  \usepackage[]{microtype}
  \UseMicrotypeSet[protrusion]{basicmath} % disable protrusion for tt fonts
}{}
\makeatletter
\@ifundefined{KOMAClassName}{% if non-KOMA class
  \IfFileExists{parskip.sty}{%
    \usepackage{parskip}
  }{% else
    \setlength{\parindent}{0pt}
    \setlength{\parskip}{6pt plus 2pt minus 1pt}}
}{% if KOMA class
  \KOMAoptions{parskip=half}}
\makeatother
\usepackage{xcolor}
\usepackage[top=20mm,left=20mm,bottom=20mm,right=20mm]{geometry}
\setlength{\emergencystretch}{3em} % prevent overfull lines
\setcounter{secnumdepth}{-\maxdimen} % remove section numbering
% Make \paragraph and \subparagraph free-standing
\makeatletter
\ifx\paragraph\undefined\else
  \let\oldparagraph\paragraph
  \renewcommand{\paragraph}{
    \@ifstar
      \xxxParagraphStar
      \xxxParagraphNoStar
  }
  \newcommand{\xxxParagraphStar}[1]{\oldparagraph*{#1}\mbox{}}
  \newcommand{\xxxParagraphNoStar}[1]{\oldparagraph{#1}\mbox{}}
\fi
\ifx\subparagraph\undefined\else
  \let\oldsubparagraph\subparagraph
  \renewcommand{\subparagraph}{
    \@ifstar
      \xxxSubParagraphStar
      \xxxSubParagraphNoStar
  }
  \newcommand{\xxxSubParagraphStar}[1]{\oldsubparagraph*{#1}\mbox{}}
  \newcommand{\xxxSubParagraphNoStar}[1]{\oldsubparagraph{#1}\mbox{}}
\fi
\makeatother

\usepackage{color}
\usepackage{fancyvrb}
\newcommand{\VerbBar}{|}
\newcommand{\VERB}{\Verb[commandchars=\\\{\}]}
\DefineVerbatimEnvironment{Highlighting}{Verbatim}{commandchars=\\\{\}}
% Add ',fontsize=\small' for more characters per line
\usepackage{framed}
\definecolor{shadecolor}{RGB}{241,243,245}
\newenvironment{Shaded}{\begin{snugshade}}{\end{snugshade}}
\newcommand{\AlertTok}[1]{\textcolor[rgb]{0.68,0.00,0.00}{#1}}
\newcommand{\AnnotationTok}[1]{\textcolor[rgb]{0.37,0.37,0.37}{#1}}
\newcommand{\AttributeTok}[1]{\textcolor[rgb]{0.40,0.45,0.13}{#1}}
\newcommand{\BaseNTok}[1]{\textcolor[rgb]{0.68,0.00,0.00}{#1}}
\newcommand{\BuiltInTok}[1]{\textcolor[rgb]{0.00,0.23,0.31}{#1}}
\newcommand{\CharTok}[1]{\textcolor[rgb]{0.13,0.47,0.30}{#1}}
\newcommand{\CommentTok}[1]{\textcolor[rgb]{0.37,0.37,0.37}{#1}}
\newcommand{\CommentVarTok}[1]{\textcolor[rgb]{0.37,0.37,0.37}{\textit{#1}}}
\newcommand{\ConstantTok}[1]{\textcolor[rgb]{0.56,0.35,0.01}{#1}}
\newcommand{\ControlFlowTok}[1]{\textcolor[rgb]{0.00,0.23,0.31}{\textbf{#1}}}
\newcommand{\DataTypeTok}[1]{\textcolor[rgb]{0.68,0.00,0.00}{#1}}
\newcommand{\DecValTok}[1]{\textcolor[rgb]{0.68,0.00,0.00}{#1}}
\newcommand{\DocumentationTok}[1]{\textcolor[rgb]{0.37,0.37,0.37}{\textit{#1}}}
\newcommand{\ErrorTok}[1]{\textcolor[rgb]{0.68,0.00,0.00}{#1}}
\newcommand{\ExtensionTok}[1]{\textcolor[rgb]{0.00,0.23,0.31}{#1}}
\newcommand{\FloatTok}[1]{\textcolor[rgb]{0.68,0.00,0.00}{#1}}
\newcommand{\FunctionTok}[1]{\textcolor[rgb]{0.28,0.35,0.67}{#1}}
\newcommand{\ImportTok}[1]{\textcolor[rgb]{0.00,0.46,0.62}{#1}}
\newcommand{\InformationTok}[1]{\textcolor[rgb]{0.37,0.37,0.37}{#1}}
\newcommand{\KeywordTok}[1]{\textcolor[rgb]{0.00,0.23,0.31}{\textbf{#1}}}
\newcommand{\NormalTok}[1]{\textcolor[rgb]{0.00,0.23,0.31}{#1}}
\newcommand{\OperatorTok}[1]{\textcolor[rgb]{0.37,0.37,0.37}{#1}}
\newcommand{\OtherTok}[1]{\textcolor[rgb]{0.00,0.23,0.31}{#1}}
\newcommand{\PreprocessorTok}[1]{\textcolor[rgb]{0.68,0.00,0.00}{#1}}
\newcommand{\RegionMarkerTok}[1]{\textcolor[rgb]{0.00,0.23,0.31}{#1}}
\newcommand{\SpecialCharTok}[1]{\textcolor[rgb]{0.37,0.37,0.37}{#1}}
\newcommand{\SpecialStringTok}[1]{\textcolor[rgb]{0.13,0.47,0.30}{#1}}
\newcommand{\StringTok}[1]{\textcolor[rgb]{0.13,0.47,0.30}{#1}}
\newcommand{\VariableTok}[1]{\textcolor[rgb]{0.07,0.07,0.07}{#1}}
\newcommand{\VerbatimStringTok}[1]{\textcolor[rgb]{0.13,0.47,0.30}{#1}}
\newcommand{\WarningTok}[1]{\textcolor[rgb]{0.37,0.37,0.37}{\textit{#1}}}

\providecommand{\tightlist}{%
  \setlength{\itemsep}{0pt}\setlength{\parskip}{0pt}}\usepackage{longtable,booktabs,array}
\usepackage{calc} % for calculating minipage widths
% Correct order of tables after \paragraph or \subparagraph
\usepackage{etoolbox}
\makeatletter
\patchcmd\longtable{\par}{\if@noskipsec\mbox{}\fi\par}{}{}
\makeatother
% Allow footnotes in longtable head/foot
\IfFileExists{footnotehyper.sty}{\usepackage{footnotehyper}}{\usepackage{footnote}}
\makesavenoteenv{longtable}
\usepackage{graphicx}
\makeatletter
\def\maxwidth{\ifdim\Gin@nat@width>\linewidth\linewidth\else\Gin@nat@width\fi}
\def\maxheight{\ifdim\Gin@nat@height>\textheight\textheight\else\Gin@nat@height\fi}
\makeatother
% Scale images if necessary, so that they will not overflow the page
% margins by default, and it is still possible to overwrite the defaults
% using explicit options in \includegraphics[width, height, ...]{}
\setkeys{Gin}{width=\maxwidth,height=\maxheight,keepaspectratio}
% Set default figure placement to htbp
\makeatletter
\def\fps@figure{htbp}
\makeatother

\KOMAoption{captions}{tableheading}
\makeatletter
\@ifpackageloaded{caption}{}{\usepackage{caption}}
\AtBeginDocument{%
\ifdefined\contentsname
  \renewcommand*\contentsname{Table of contents}
\else
  \newcommand\contentsname{Table of contents}
\fi
\ifdefined\listfigurename
  \renewcommand*\listfigurename{List of Figures}
\else
  \newcommand\listfigurename{List of Figures}
\fi
\ifdefined\listtablename
  \renewcommand*\listtablename{List of Tables}
\else
  \newcommand\listtablename{List of Tables}
\fi
\ifdefined\figurename
  \renewcommand*\figurename{Figure}
\else
  \newcommand\figurename{Figure}
\fi
\ifdefined\tablename
  \renewcommand*\tablename{Table}
\else
  \newcommand\tablename{Table}
\fi
}
\@ifpackageloaded{float}{}{\usepackage{float}}
\floatstyle{ruled}
\@ifundefined{c@chapter}{\newfloat{codelisting}{h}{lop}}{\newfloat{codelisting}{h}{lop}[chapter]}
\floatname{codelisting}{Listing}
\newcommand*\listoflistings{\listof{codelisting}{List of Listings}}
\makeatother
\makeatletter
\makeatother
\makeatletter
\@ifpackageloaded{caption}{}{\usepackage{caption}}
\@ifpackageloaded{subcaption}{}{\usepackage{subcaption}}
\makeatother

\ifLuaTeX
  \usepackage{selnolig}  % disable illegal ligatures
\fi
\usepackage{bookmark}

\IfFileExists{xurl.sty}{\usepackage{xurl}}{} % add URL line breaks if available
\urlstyle{same} % disable monospaced font for URLs
\hypersetup{
  pdftitle={Marginalizing shocks},
  colorlinks=true,
  linkcolor={blue},
  filecolor={Maroon},
  citecolor={Blue},
  urlcolor={Blue},
  pdfcreator={LaTeX via pandoc}}


\title{Marginalizing shocks}
\author{}
\date{}

\begin{document}
\maketitle


\emph{Cover that shocks which I ought not to see.}

\linespread{0.5}

\linespread{2}

\subsection{Model reparameterization}\label{model-reparameterization}

We marginalize the shocks in the mixture model! Going from:

\linespread{0.5}

\begin{Shaded}
\begin{Highlighting}[]
\CommentTok{\# Non{-}shocky}
\KeywordTok{target +=}\NormalTok{ normal\_lpdf(y[n] | mu[n] + delta\_sck[n], sigma)}
\KeywordTok{target +=}\NormalTok{ normal\_lpdf(delta\_sck[n] | }\DecValTok{0}\NormalTok{, inner\_tau\_sck)}

\CommentTok{\# Shocky}
\KeywordTok{target +=}\NormalTok{ normal\_lpdf(y[n] | mu[n] + delta\_sck[n], sigma)}
\KeywordTok{target +=}\NormalTok{ normal\_lpdf(delta\_sck[n] | }\DecValTok{0}\NormalTok{, outer\_tau\_sck)}
\end{Highlighting}
\end{Shaded}

\linespread{2}

to a log-likelihood where shocks are marginalized out:

\linespread{0.5}

\begin{Shaded}
\begin{Highlighting}[]
\CommentTok{\# Non{-}shocky}
\KeywordTok{target +=}\NormalTok{ normal\_lpdf(y[n] | mu[n], eta)}

\CommentTok{\# Shocky}
\KeywordTok{target +=}\NormalTok{ normal\_lpdf( y[n] | mu[n], eta * sqrt\{}\DecValTok{1}\NormalTok{ + nu / eta\^{}}\DecValTok{2}\NormalTok{\})}
\end{Highlighting}
\end{Shaded}

\linespread{2}

where: \[\eta = \sqrt{ \sigma^2 + {\tau_\text{inner}^\text{shock}}^2 }\]
\[\nu = {\tau_\text{outer}^\text{shock}}^2-{\tau_\text{inner}^\text{shock}}^2\]
In Stan, we now have in the model block:

\linespread{0.5}

\begin{Shaded}
\begin{Highlighting}[]
\FunctionTok{writeLines}\NormalTok{(}
  \FunctionTok{readLines}\NormalTok{(}\FunctionTok{file.path}\NormalTok{(wd, }\StringTok{\textquotesingle{}model\textquotesingle{}}\NormalTok{, }\StringTok{\textquotesingle{}stan\textquotesingle{}}\NormalTok{, }
                      \StringTok{\textquotesingle{}model9\_marginalshocks.stan\textquotesingle{}}\NormalTok{))[}\DecValTok{230}\SpecialCharTok{:}\DecValTok{315}\NormalTok{])}
\end{Highlighting}
\end{Shaded}

\begin{verbatim}
model {
  
  alpha ~ normal(0, 0.69); // 0.2 mm <~ exp(alpha) * 1 mm <~ 5 mm
  
  beta_gdd ~ normal(0, log(1.8) / 2.57); // -log(1.8) <~ beta_gsl <~ log(1.8)
  beta_sm ~ normal(0, log(1.8) / 2.57); // -log(1.8) <~ beta_gsl <~ log(1.8)
  beta_vpd ~ normal(0, log(1.8) / 2.57); // -log(1.8) <~ beta_gsl <~ log(1.8)
  
  rho_sp ~ lognormal(2.65, 0.135); // 10 <~ rho <~ 20
  
  gamma_sp ~ normal(0, log(10) / 2.57); // 0 <~ gamma <~ log(10)
  
  for (s in 1:N_stands)
    f_tilde_sh[s] ~ normal(0, 1);
  rho_sh ~ lognormal(1.7, 0.26);       // 3 <~ rho_sh <~ 10
  gamma_sh ~ normal(0, log(3) / 2.57); // 0 <~ gamma_sh <~ log(3)
  
  eta ~ gamma(6, 43.5);
  nu ~ gamma(0.5, 0.033); 
  phi_sck ~ beta(2, 20); 
  omega_conc_sck ~ beta(230, 14); // 0.9 <~ omega_conc_sck <~ 0.97 (most trees, but not ALL trees)
  omega_nonconc_sck ~ beta(1, 20); // 0 <~ omega_conc_sck <~ 0.15 (should be rare, but... who knows?)
  
  vector[N] mu;
  for (t in 1:N_trees) {
    array[N_years[t]] int tree_idxs
    = linspaced_int_array(N_years[t],
                          tree_start_idxs[t],
                          tree_end_idxs[t]);

    array[N_years[t]] int all_years_idxs_tree = all_years_idxs[tree_idxs];
    array[N_years[t]] real years_tree = years[tree_idxs];

    int stand_idx = stand_idxs[t];
    int species_idx = species_idxs[t];

    vector[N_years[t]] gdd_obs_tree = gdd_obs[tree_idxs];
    vector[N_years[t]] sm_obs_tree = sm_obs[tree_idxs];
    vector[N_years[t]] vpd_obs_tree = vpd_obs[tree_idxs];

    mu[tree_idxs] =  alpha
    + beta_gdd[species_idx] * (gdd_obs_tree - gdd0)
    + beta_sm[species_idx] * (sm_obs_tree - sm0)
    + beta_vpd[species_idx] * (vpd_obs_tree - vpd0)
    + f_sh[stand_idx, all_years_idxs_tree];
    
    f_tilde[tree_idxs] ~ normal(0,1);
  }
  
  // Observational model with marginalized shocks
  for (s in 1:N_stands) {
    
    array[N_stand_trees[s]] int stand_trees_idxs = linspaced_int_array(N_stand_trees[s], 
      stand_trees_start_idxs[s], stand_trees_end_idxs[s]);
    vector[N_stand_years[s]] log_p0 = rep_vector(0, N_stand_years[s]);
    vector[N_stand_years[s]] log_p1 = rep_vector(0, N_stand_years[s]);
    
    for(ts in 1:N_stand_trees[s]){
      int t = stand_trees_idxs[ts];
      array[N_years[t]] int tree_idxs = linspaced_int_array(N_years[t], tree_start_idxs[t], tree_end_idxs[t]);
      array[N_years[t]] int all_years_idxs_tree = all_years_idxs[tree_idxs];
      
      for(y in 1:N_years[t]) {
        int ys = all_years_idxs_tree[y]-stand_start_years_idxs[s]+1;
        
        log_p0[ys] += log_mix(omega_nonconc_sck,
                          normal_lpdf(log_rw_obs[tree_idxs[y]] | mu[tree_idxs[y]] 
                          + f[tree_idxs[y]], eta * sqrt(1 + nu / eta^2)),
                          normal_lpdf(log_rw_obs[tree_idxs[y]] | mu[tree_idxs[y]] 
                          + f[tree_idxs[y]], eta));
        log_p1[ys] += log_mix(omega_conc_sck,
                          normal_lpdf(log_rw_obs[tree_idxs[y]] | mu[tree_idxs[y]] 
                          + f[tree_idxs[y]], eta * sqrt(1 + nu / eta^2)),
                          normal_lpdf(log_rw_obs[tree_idxs[y]] | mu[tree_idxs[y]] 
                          + f[tree_idxs[y]], eta));
        
      }
    }
    
    for(y in 1:N_stand_years[s]) {
      target += log_mix(phi_sck[s], log_p1[y], log_p0[y]);
    }
  }
  
}
\end{verbatim}

\linespread{2}

\subsection{But\ldots{} my shocks\ldots{}}\label{but-my-shocks}

Our large shocks are likely tied to interesting ecological stuff\ldots{}
But by marginalizing them, we don't have inferences on them anymore.\\
It could be doable to ``\emph{to marginalize out eta in the ``small
shock'' branch of the mixture model but leave eta explicit in the
``large shock'' branch of the mixture model}'' (Mike).\\
Since I'm not totally sure how to do that, I decided to sample the
\(\delta_\text{shock}\) afterwards, to recover their posterior
distributions in the generated quantities block. This is an ongoing
project, so for now I recover only the shock state (\(z_t\)).

\linespread{0.5}

\begin{Shaded}
\begin{Highlighting}[]
\FunctionTok{writeLines}\NormalTok{(}
  \FunctionTok{readLines}\NormalTok{(}\FunctionTok{file.path}\NormalTok{(wd, }\StringTok{\textquotesingle{}model\textquotesingle{}}\NormalTok{, }\StringTok{\textquotesingle{}stan\textquotesingle{}}\NormalTok{, }
                      \StringTok{\textquotesingle{}model9\_marginalshocks.stan\textquotesingle{}}\NormalTok{))[}\DecValTok{351}\SpecialCharTok{:}\DecValTok{405}\NormalTok{])}
\end{Highlighting}
\end{Shaded}

\begin{verbatim}
generated quantities {

  array[N] int sck_state; // latent state, zt = 0 or zt = 1
  array[N] real log_rw_pred;

  for (t in 1:N_trees) {
    
    array[N_years[t]] int tree_idxs
    = linspaced_int_array(N_years[t],
                          tree_start_idxs[t],
                          tree_end_idxs[t]);

    array[N_years[t]] int all_years_idxs_tree = all_years_idxs[tree_idxs];
    array[N_years[t]] real years_tree = years[tree_idxs];

    int stand_idx = stand_idxs[t];
    int species_idx = species_idxs[t];

    vector[N_years[t]] gdd_obs_tree = gdd_obs[tree_idxs];
    vector[N_years[t]] sm_obs_tree = sm_obs[tree_idxs];
    vector[N_years[t]] vpd_obs_tree = vpd_obs[tree_idxs];

    vector[N_years[t]] mu =  alpha
    + beta_gdd[species_idx] * (gdd_obs_tree - gdd0)
    + beta_sm[species_idx] * (sm_obs_tree - sm0)
    + beta_vpd[species_idx] * (vpd_obs_tree - vpd0)
    + f_sh[stand_idx, all_years_idxs_tree];
    
    // mixture weight for shock
    real mw_shock = phi_sck[stand_idx]*omega_conc_sck + (1-phi_sck[stand_idx])*omega_nonconc_sck;
    
    for(y in 1:N_years[t]){
      
      // mixture components
      real log_pshock = log(mw_shock) + normal_lpdf(log_rw_obs[tree_idxs[y]] | mu[y] + f[tree_idxs[y]], eta * sqrt(1 + nu / eta^2));
      real log_pshock_plus_pnonshock = log_mix(mw_shock, 
        normal_lpdf(log_rw_obs[tree_idxs[y]] | mu[y] + f[tree_idxs[y]], eta * sqrt(1 + nu / eta^2)),
        normal_lpdf(log_rw_obs[tree_idxs[y]] | mu[y] + f[tree_idxs[y]], eta));
      
      // probability of shock state
      real lambda_shock = exp(log_pshock - log_pshock_plus_pnonshock);
      
      sck_state[tree_idxs[y]] = bernoulli_rng(lambda_shock); // or something like categorical_rng(lambda_shock);?

      
      if(sck_state[tree_idxs[y]] == 1) {
        log_rw_pred[tree_idxs[y]] = normal_rng(mu[y] + f[tree_idxs[y]], eta * sqrt(1 + nu / eta^2));
      }else{
        log_rw_pred[tree_idxs[y]] = normal_rng(mu[y] + f[tree_idxs[y]], eta);
      }

      
    }
  }
}
\end{verbatim}

\linespread{2}

\subsection{Test with a small dataset}\label{test-with-a-small-dataset}

Let's run this new model on a very small dataset, 21 trees and 1 site.\\
The good thing here is that I don't need to tune the (non-)centered
parametrization anymore.

\linespread{0.5}

\begin{Shaded}
\begin{Highlighting}[]
\NormalTok{data }\OtherTok{\textless{}{-}} \FunctionTok{readRDS}\NormalTok{(}\AttributeTok{file =} \FunctionTok{file.path}\NormalTok{(wd, }\StringTok{\textquotesingle{}output\textquotesingle{}}\NormalTok{, }\StringTok{\textquotesingle{}model\textquotesingle{}}\NormalTok{, }\StringTok{\textquotesingle{}data\_15nov2025\_az592.rds\textquotesingle{}}\NormalTok{))}
\NormalTok{data}\SpecialCharTok{$}\NormalTok{N\_stand\_trees }\OtherTok{\textless{}{-}} \FunctionTok{array}\NormalTok{(data}\SpecialCharTok{$}\NormalTok{N\_stand\_trees, }\AttributeTok{dim =} \DecValTok{1}\NormalTok{)}
\NormalTok{data}\SpecialCharTok{$}\NormalTok{N\_stand\_years }\OtherTok{\textless{}{-}} \FunctionTok{array}\NormalTok{(data}\SpecialCharTok{$}\NormalTok{N\_stand\_years, }\AttributeTok{dim =} \DecValTok{1}\NormalTok{)}
\NormalTok{data}\SpecialCharTok{$}\NormalTok{stand\_start\_years\_idxs }\OtherTok{\textless{}{-}} \FunctionTok{array}\NormalTok{(data}\SpecialCharTok{$}\NormalTok{stand\_start\_years\_idxs, }\AttributeTok{dim =} \DecValTok{1}\NormalTok{)}
\NormalTok{data}\SpecialCharTok{$}\NormalTok{stand\_trees\_end\_idxs }\OtherTok{\textless{}{-}} \FunctionTok{array}\NormalTok{(data}\SpecialCharTok{$}\NormalTok{stand\_trees\_end\_idxs, }\AttributeTok{dim =} \DecValTok{1}\NormalTok{)}
\NormalTok{data}\SpecialCharTok{$}\NormalTok{stand\_trees\_start\_idxs }\OtherTok{\textless{}{-}} \FunctionTok{array}\NormalTok{(data}\SpecialCharTok{$}\NormalTok{stand\_trees\_start\_idxs, }\AttributeTok{dim =} \DecValTok{1}\NormalTok{)}

\ControlFlowTok{if}\NormalTok{(}\ConstantTok{FALSE}\NormalTok{)\{}
\NormalTok{  fit }\OtherTok{\textless{}{-}} \FunctionTok{stan}\NormalTok{(}\AttributeTok{file=}\FunctionTok{file.path}\NormalTok{(wd, }\StringTok{\textquotesingle{}model\textquotesingle{}}\NormalTok{, }\StringTok{\textquotesingle{}stan\textquotesingle{}}\NormalTok{, }\StringTok{\textquotesingle{}model9\_marginalshocks.stan\textquotesingle{}}\NormalTok{),}
            \AttributeTok{data=}\NormalTok{data, }\AttributeTok{seed=}\DecValTok{5838293}\NormalTok{, }\AttributeTok{chains =} \DecValTok{4}\NormalTok{,}
            \AttributeTok{warmup=}\DecValTok{1000}\NormalTok{, }\AttributeTok{iter=}\DecValTok{2024}\NormalTok{, }\AttributeTok{refresh=}\DecValTok{10}\NormalTok{)}
  \FunctionTok{saveRDS}\NormalTok{(fit, }\FunctionTok{file.path}\NormalTok{(wd, }\StringTok{\textquotesingle{}model/shocks/output/marginal\textquotesingle{}}\NormalTok{, }\StringTok{\textquotesingle{}fit.rds\textquotesingle{}}\NormalTok{))}
\NormalTok{\}}\ControlFlowTok{else}\NormalTok{\{}
\NormalTok{  fit }\OtherTok{\textless{}{-}} \FunctionTok{readRDS}\NormalTok{(}\FunctionTok{file.path}\NormalTok{(wd, }\StringTok{\textquotesingle{}model/shocks/output/marginal\textquotesingle{}}\NormalTok{, }\StringTok{\textquotesingle{}fit.rds\textquotesingle{}}\NormalTok{))}
\NormalTok{\}}
\end{Highlighting}
\end{Shaded}

\linespread{2}

Let's look at the diagnostics:

\linespread{0.5}

\begin{Shaded}
\begin{Highlighting}[]
\NormalTok{diagnostics }\OtherTok{\textless{}{-}}\NormalTok{ util}\SpecialCharTok{$}\FunctionTok{extract\_hmc\_diagnostics}\NormalTok{(fit)}
\NormalTok{util}\SpecialCharTok{$}\FunctionTok{check\_all\_hmc\_diagnostics}\NormalTok{(diagnostics)}
\end{Highlighting}
\end{Shaded}

\begin{verbatim}
  All Hamiltonian Monte Carlo diagnostics are consistent with reliable
Markov chain Monte Carlo.
\end{verbatim}

\begin{Shaded}
\begin{Highlighting}[]
\NormalTok{samples }\OtherTok{\textless{}{-}}\NormalTok{ util}\SpecialCharTok{$}\FunctionTok{extract\_expectand\_vals}\NormalTok{(fit)}
\NormalTok{base\_samples }\OtherTok{\textless{}{-}}\NormalTok{ util}\SpecialCharTok{$}\FunctionTok{filter\_expectands}\NormalTok{(samples,}
                                       \FunctionTok{c}\NormalTok{(}\StringTok{\textquotesingle{}alpha\textquotesingle{}}\NormalTok{, }
                                         \StringTok{\textquotesingle{}beta\_gdd\textquotesingle{}}\NormalTok{, }\StringTok{\textquotesingle{}beta\_sm\textquotesingle{}}\NormalTok{, }\StringTok{\textquotesingle{}beta\_vpd\textquotesingle{}}\NormalTok{,}
                                         \StringTok{\textquotesingle{}rho\_sh\textquotesingle{}}\NormalTok{, }\StringTok{\textquotesingle{}gamma\_sh\textquotesingle{}}\NormalTok{,}
                                         \StringTok{\textquotesingle{}rho\_sp\textquotesingle{}}\NormalTok{, }\StringTok{\textquotesingle{}gamma\_sp\textquotesingle{}}\NormalTok{,}
                                         \StringTok{\textquotesingle{}eta\textquotesingle{}}\NormalTok{, }\StringTok{\textquotesingle{}nu\textquotesingle{}}\NormalTok{,}
                                         \StringTok{\textquotesingle{}phi\_sck\textquotesingle{}}\NormalTok{, }\StringTok{\textquotesingle{}omega\_conc\_sck\textquotesingle{}}\NormalTok{, }\StringTok{\textquotesingle{}omega\_nonconc\_sck\textquotesingle{}}\NormalTok{),}
                                       \AttributeTok{check\_arrays =} \ConstantTok{TRUE}\NormalTok{)}
\NormalTok{util}\SpecialCharTok{$}\FunctionTok{check\_all\_expectand\_diagnostics}\NormalTok{(base\_samples)}
\end{Highlighting}
\end{Shaded}

\begin{verbatim}
All expectands checked appear to be behaving well enough for reliable
Markov chain Monte Carlo estimation.
\end{verbatim}

\linespread{2}

Clean! Let's continue.\\
These are the latent shock states:

\linespread{0.5}

\begin{Shaded}
\begin{Highlighting}[]
\FunctionTok{set.seed}\NormalTok{(}\DecValTok{123456}\NormalTok{)}
\NormalTok{random\_trees }\OtherTok{\textless{}{-}} \FunctionTok{sample}\NormalTok{(}\DecValTok{1}\SpecialCharTok{:}\NormalTok{data}\SpecialCharTok{$}\NormalTok{N\_trees, }\DecValTok{4}\NormalTok{)}

\FunctionTok{par}\NormalTok{(}\AttributeTok{mfrow =} \FunctionTok{c}\NormalTok{(}\DecValTok{2}\NormalTok{,}\DecValTok{2}\NormalTok{), }\AttributeTok{mar =} \FunctionTok{c}\NormalTok{(}\DecValTok{2}\NormalTok{, }\DecValTok{5}\NormalTok{, }\DecValTok{3}\NormalTok{, }\DecValTok{1}\NormalTok{))}
\ControlFlowTok{for}\NormalTok{(t }\ControlFlowTok{in}\NormalTok{ random\_trees)\{}
\NormalTok{  idxs\_t }\OtherTok{\textless{}{-}}\NormalTok{ data}\SpecialCharTok{$}\NormalTok{tree\_start\_idxs[t]}\SpecialCharTok{:}\NormalTok{data}\SpecialCharTok{$}\NormalTok{tree\_end\_idxs[t]}
\NormalTok{  names }\OtherTok{\textless{}{-}} \FunctionTok{sapply}\NormalTok{(idxs\_t,}
                  \ControlFlowTok{function}\NormalTok{(x) }\FunctionTok{paste0}\NormalTok{(}\StringTok{\textquotesingle{}sck\_state[\textquotesingle{}}\NormalTok{, x, }\StringTok{\textquotesingle{}]\textquotesingle{}}\NormalTok{))}
\NormalTok{  util}\SpecialCharTok{$}\FunctionTok{plot\_conn\_pushforward\_quantiles}\NormalTok{(samples, names, data}\SpecialCharTok{$}\NormalTok{years[idxs\_t],}
                                       \AttributeTok{xlab=}\StringTok{""}\NormalTok{, }\AttributeTok{ylab=}\FunctionTok{paste0}\NormalTok{(}\StringTok{"Shock state {-} tree "}\NormalTok{,t), }
                                       \AttributeTok{display\_ylim=}\FunctionTok{c}\NormalTok{(}\SpecialCharTok{{-}}\FloatTok{0.1}\NormalTok{, }\FloatTok{1.1}\NormalTok{), }\AttributeTok{display\_xlim =} \FunctionTok{c}\NormalTok{(}\DecValTok{1980}\NormalTok{, }\DecValTok{2010}\NormalTok{))}
\NormalTok{\}}
\end{Highlighting}
\end{Shaded}

\includegraphics{marginal_shocks_files/figure-pdf/unnamed-chunk-8-1.pdf}

\linespread{2}

\subsection{Next step}\label{next-step}

Is it possible to recover the posterior of
\(\delta_\text{shock}(k) ~|~ y(k)\), without having
\({\tau_\text{inner}^\text{shock}}\) or
\({\tau_\text{outer}^\text{shock}}\)?\\
Maybe a first step could be to fix:
\({\tau_\text{inner}^\text{shock}} = 0\)?

Because for now my predictions look like:

\linespread{0.5}

\begin{Shaded}
\begin{Highlighting}[]
\FunctionTok{par}\NormalTok{(}\AttributeTok{mfrow =} \FunctionTok{c}\NormalTok{(}\DecValTok{2}\NormalTok{,}\DecValTok{2}\NormalTok{), }\AttributeTok{mar =} \FunctionTok{c}\NormalTok{(}\DecValTok{2}\NormalTok{, }\DecValTok{5}\NormalTok{, }\DecValTok{3}\NormalTok{, }\DecValTok{1}\NormalTok{))}
\ControlFlowTok{for}\NormalTok{(t }\ControlFlowTok{in}\NormalTok{ random\_trees)\{}
\NormalTok{  idxs\_t }\OtherTok{\textless{}{-}}\NormalTok{ data}\SpecialCharTok{$}\NormalTok{tree\_start\_idxs[t]}\SpecialCharTok{:}\NormalTok{data}\SpecialCharTok{$}\NormalTok{tree\_end\_idxs[t]}
\NormalTok{  names }\OtherTok{\textless{}{-}} \FunctionTok{sapply}\NormalTok{(idxs\_t,}
                  \ControlFlowTok{function}\NormalTok{(x) }\FunctionTok{paste0}\NormalTok{(}\StringTok{\textquotesingle{}log\_rw\_pred[\textquotesingle{}}\NormalTok{, x, }\StringTok{\textquotesingle{}]\textquotesingle{}}\NormalTok{))}
\NormalTok{  util}\SpecialCharTok{$}\FunctionTok{plot\_conn\_pushforward\_quantiles}\NormalTok{(samples, names, data}\SpecialCharTok{$}\NormalTok{years[idxs\_t],}
                                       \AttributeTok{xlab=}\StringTok{""}\NormalTok{, }\AttributeTok{ylab=}\FunctionTok{paste0}\NormalTok{(}\StringTok{"Shock state {-} tree "}\NormalTok{,t), }
                                       \AttributeTok{display\_ylim=}\FunctionTok{c}\NormalTok{(}\SpecialCharTok{{-}}\DecValTok{10}\NormalTok{, }\DecValTok{5}\NormalTok{), }\AttributeTok{display\_xlim =} \FunctionTok{c}\NormalTok{(}\DecValTok{1980}\NormalTok{, }\DecValTok{2010}\NormalTok{))}
  \FunctionTok{points}\NormalTok{(data}\SpecialCharTok{$}\NormalTok{years[idxs\_t], data}\SpecialCharTok{$}\NormalTok{log\_rw\_obs[idxs\_t], }\AttributeTok{pch=}\DecValTok{16}\NormalTok{, }\AttributeTok{cex=}\FloatTok{0.7}\NormalTok{, }\AttributeTok{col=}\StringTok{"white"}\NormalTok{)}
  \FunctionTok{points}\NormalTok{(data}\SpecialCharTok{$}\NormalTok{years[idxs\_t], data}\SpecialCharTok{$}\NormalTok{log\_rw\_obs[idxs\_t], }\AttributeTok{pch=}\DecValTok{16}\NormalTok{, }\AttributeTok{cex=}\FloatTok{0.4}\NormalTok{, }\AttributeTok{col=}\StringTok{"black"}\NormalTok{)}
\NormalTok{\}}
\end{Highlighting}
\end{Shaded}

\includegraphics{marginal_shocks_files/figure-pdf/unnamed-chunk-9-1.pdf}

\linespread{2}

\subsection{Assuming that tau\_inner =
0}\label{assuming-that-tau_inner-0}

If I assume that \({\tau_\text{inner}^\text{shock}} = 0\), I can
reconstruct the large shocks using the normal-normal conjugancy
(\emph{merci} Google):
\[\delta_\text{shock}(k) ~|~ y_k \sim \text{normal} \left( ({\tau_\text{outer}^\text{shock}}^2 / ({\tau_\text{outer}^\text{shock}}^2 + \sigma^2)) \times y_k ~,~ \sqrt{({\tau_\text{outer}^\text{shock}}^2 \times \sigma^2) / ({\tau_\text{outer}^\text{shock}}^2 + \sigma^2)}\right)\]

In the generated quantities block, we can then have:

\linespread{0.5}

\begin{Shaded}
\begin{Highlighting}[]
\FunctionTok{writeLines}\NormalTok{(}
  \FunctionTok{readLines}\NormalTok{(}\FunctionTok{file.path}\NormalTok{(wd, }\StringTok{\textquotesingle{}model\textquotesingle{}}\NormalTok{, }\StringTok{\textquotesingle{}stan\textquotesingle{}}\NormalTok{, }
                      \StringTok{\textquotesingle{}model9\_marginalshocks\_zerotauinner.stan\textquotesingle{}}\NormalTok{))[}\DecValTok{395}\SpecialCharTok{:}\DecValTok{405}\NormalTok{])}
\end{Highlighting}
\end{Shaded}

\begin{verbatim}

      
      if(sck_state[tree_idxs[y]] == 1) {
        // we can reconstruct shock posterior using the normal-normal conjugancy 
        real conjugate_mean = (outer_tau_sck^2 / (outer_tau_sck^2 + sigma^2)) * log_rw_obs[tree_idxs[y]];
        real conjugate_sd   = sqrt((outer_tau_sck^2 * sigma^2) / (outer_tau_sck^2 + sigma^2));
        delta_sck[tree_idxs[y]] = normal_rng(conjugate_mean, conjugate_sd);
        log_rw_pred[tree_idxs[y]] = normal_rng(mu[y] + f[tree_idxs[y]] + delta_sck[tree_idxs[y]], sigma);
      }else{
        log_rw_pred[tree_idxs[y]] = normal_rng(mu[y] + f[tree_idxs[y]], sigma);
      }
\end{verbatim}

\linespread{2}

\linespread{0.5}

\begin{Shaded}
\begin{Highlighting}[]
\ControlFlowTok{if}\NormalTok{(}\ConstantTok{FALSE}\NormalTok{)\{}
\NormalTok{  fit2 }\OtherTok{\textless{}{-}} \FunctionTok{stan}\NormalTok{(}\AttributeTok{file=}\FunctionTok{file.path}\NormalTok{(wd, }\StringTok{\textquotesingle{}model\textquotesingle{}}\NormalTok{, }\StringTok{\textquotesingle{}stan\textquotesingle{}}\NormalTok{, }\StringTok{\textquotesingle{}model9\_marginalshocks\_zerotauinner.stan\textquotesingle{}}\NormalTok{),}
            \AttributeTok{data=}\NormalTok{data, }\AttributeTok{seed=}\DecValTok{5838293}\NormalTok{, }\AttributeTok{chains =} \DecValTok{4}\NormalTok{,}
            \AttributeTok{warmup=}\DecValTok{1000}\NormalTok{, }\AttributeTok{iter=}\DecValTok{2024}\NormalTok{, }\AttributeTok{refresh=}\DecValTok{10}\NormalTok{)}
  \FunctionTok{saveRDS}\NormalTok{(fit2, }\FunctionTok{file.path}\NormalTok{(wd, }\StringTok{\textquotesingle{}model/shocks/output/marginal\textquotesingle{}}\NormalTok{, }\StringTok{\textquotesingle{}fit2.rds\textquotesingle{}}\NormalTok{))}
\NormalTok{\}}\ControlFlowTok{else}\NormalTok{\{}
\NormalTok{  fit2 }\OtherTok{\textless{}{-}} \FunctionTok{readRDS}\NormalTok{(}\FunctionTok{file.path}\NormalTok{(wd, }\StringTok{\textquotesingle{}model/shocks/output/marginal\textquotesingle{}}\NormalTok{, }\StringTok{\textquotesingle{}fit2.rds\textquotesingle{}}\NormalTok{))}
\NormalTok{\}}
\end{Highlighting}
\end{Shaded}

\linespread{2}

Diagnostics are cleaned:

\linespread{0.5}

\begin{Shaded}
\begin{Highlighting}[]
\NormalTok{diagnostics }\OtherTok{\textless{}{-}}\NormalTok{ util}\SpecialCharTok{$}\FunctionTok{extract\_hmc\_diagnostics}\NormalTok{(fit2)}
\NormalTok{util}\SpecialCharTok{$}\FunctionTok{check\_all\_hmc\_diagnostics}\NormalTok{(diagnostics)}
\end{Highlighting}
\end{Shaded}

\begin{verbatim}
  All Hamiltonian Monte Carlo diagnostics are consistent with reliable
Markov chain Monte Carlo.
\end{verbatim}

\begin{Shaded}
\begin{Highlighting}[]
\NormalTok{samples }\OtherTok{\textless{}{-}}\NormalTok{ util}\SpecialCharTok{$}\FunctionTok{extract\_expectand\_vals}\NormalTok{(fit2)}
\NormalTok{base\_samples }\OtherTok{\textless{}{-}}\NormalTok{ util}\SpecialCharTok{$}\FunctionTok{filter\_expectands}\NormalTok{(samples,}
                                       \FunctionTok{c}\NormalTok{(}\StringTok{\textquotesingle{}alpha\textquotesingle{}}\NormalTok{, }
                                         \StringTok{\textquotesingle{}beta\_gdd\textquotesingle{}}\NormalTok{, }\StringTok{\textquotesingle{}beta\_sm\textquotesingle{}}\NormalTok{, }\StringTok{\textquotesingle{}beta\_vpd\textquotesingle{}}\NormalTok{,}
                                         \StringTok{\textquotesingle{}rho\_sh\textquotesingle{}}\NormalTok{, }\StringTok{\textquotesingle{}gamma\_sh\textquotesingle{}}\NormalTok{,}
                                         \StringTok{\textquotesingle{}rho\_sp\textquotesingle{}}\NormalTok{, }\StringTok{\textquotesingle{}gamma\_sp\textquotesingle{}}\NormalTok{,}
                                         \StringTok{\textquotesingle{}sigma\textquotesingle{}}\NormalTok{, }\StringTok{\textquotesingle{}outer\_tau\_sck\textquotesingle{}}\NormalTok{,}
                                         \StringTok{\textquotesingle{}phi\_sck\textquotesingle{}}\NormalTok{, }\StringTok{\textquotesingle{}omega\_conc\_sck\textquotesingle{}}\NormalTok{, }\StringTok{\textquotesingle{}omega\_nonconc\_sck\textquotesingle{}}\NormalTok{),}
                                       \AttributeTok{check\_arrays =} \ConstantTok{TRUE}\NormalTok{)}
\NormalTok{util}\SpecialCharTok{$}\FunctionTok{check\_all\_expectand\_diagnostics}\NormalTok{(base\_samples)}
\end{Highlighting}
\end{Shaded}

\begin{verbatim}
All expectands checked appear to be behaving well enough for reliable
Markov chain Monte Carlo estimation.
\end{verbatim}

\linespread{2}

We can look at the latent shocks:

\linespread{0.5}

\begin{Shaded}
\begin{Highlighting}[]
\FunctionTok{par}\NormalTok{(}\AttributeTok{mfrow =} \FunctionTok{c}\NormalTok{(}\DecValTok{2}\NormalTok{,}\DecValTok{2}\NormalTok{), }\AttributeTok{mar =} \FunctionTok{c}\NormalTok{(}\DecValTok{2}\NormalTok{, }\DecValTok{5}\NormalTok{, }\DecValTok{3}\NormalTok{, }\DecValTok{1}\NormalTok{))}
\ControlFlowTok{for}\NormalTok{(t }\ControlFlowTok{in}\NormalTok{ random\_trees)\{}
\NormalTok{  idxs\_t }\OtherTok{\textless{}{-}}\NormalTok{ data}\SpecialCharTok{$}\NormalTok{tree\_start\_idxs[t]}\SpecialCharTok{:}\NormalTok{data}\SpecialCharTok{$}\NormalTok{tree\_end\_idxs[t]}
\NormalTok{  names }\OtherTok{\textless{}{-}} \FunctionTok{sapply}\NormalTok{(idxs\_t,}
                  \ControlFlowTok{function}\NormalTok{(x) }\FunctionTok{paste0}\NormalTok{(}\StringTok{\textquotesingle{}delta\_sck[\textquotesingle{}}\NormalTok{, x, }\StringTok{\textquotesingle{}]\textquotesingle{}}\NormalTok{))}
\NormalTok{  util}\SpecialCharTok{$}\FunctionTok{plot\_conn\_pushforward\_quantiles}\NormalTok{(samples, names, data}\SpecialCharTok{$}\NormalTok{years[idxs\_t],}
                                       \AttributeTok{xlab=}\StringTok{""}\NormalTok{, }\AttributeTok{ylab=}\FunctionTok{paste0}\NormalTok{(}\StringTok{"Shock state {-} tree "}\NormalTok{,t), }
                                       \AttributeTok{display\_ylim=}\FunctionTok{c}\NormalTok{(}\SpecialCharTok{{-}}\DecValTok{10}\NormalTok{, }\FloatTok{0.1}\NormalTok{), }\AttributeTok{display\_xlim =} \FunctionTok{c}\NormalTok{(}\DecValTok{1980}\NormalTok{, }\DecValTok{2010}\NormalTok{))}
\NormalTok{\}}
\end{Highlighting}
\end{Shaded}

\includegraphics{marginal_shocks_files/figure-pdf/unnamed-chunk-13-1.pdf}

\linespread{2}

And at the predictions:

\linespread{0.5}

\begin{Shaded}
\begin{Highlighting}[]
\FunctionTok{par}\NormalTok{(}\AttributeTok{mfrow =} \FunctionTok{c}\NormalTok{(}\DecValTok{2}\NormalTok{,}\DecValTok{2}\NormalTok{), }\AttributeTok{mar =} \FunctionTok{c}\NormalTok{(}\DecValTok{2}\NormalTok{, }\DecValTok{5}\NormalTok{, }\DecValTok{3}\NormalTok{, }\DecValTok{1}\NormalTok{))}
\ControlFlowTok{for}\NormalTok{(t }\ControlFlowTok{in}\NormalTok{ random\_trees)\{}
\NormalTok{  idxs\_t }\OtherTok{\textless{}{-}}\NormalTok{ data}\SpecialCharTok{$}\NormalTok{tree\_start\_idxs[t]}\SpecialCharTok{:}\NormalTok{data}\SpecialCharTok{$}\NormalTok{tree\_end\_idxs[t]}
\NormalTok{  names }\OtherTok{\textless{}{-}} \FunctionTok{sapply}\NormalTok{(idxs\_t,}
                  \ControlFlowTok{function}\NormalTok{(x) }\FunctionTok{paste0}\NormalTok{(}\StringTok{\textquotesingle{}log\_rw\_pred[\textquotesingle{}}\NormalTok{, x, }\StringTok{\textquotesingle{}]\textquotesingle{}}\NormalTok{))}
\NormalTok{  util}\SpecialCharTok{$}\FunctionTok{plot\_conn\_pushforward\_quantiles}\NormalTok{(samples, names, data}\SpecialCharTok{$}\NormalTok{years[idxs\_t],}
                                       \AttributeTok{xlab=}\StringTok{""}\NormalTok{, }\AttributeTok{ylab=}\FunctionTok{paste0}\NormalTok{(}\StringTok{"log(rw) {-} tree "}\NormalTok{,t), }
                                       \AttributeTok{display\_ylim=}\FunctionTok{c}\NormalTok{(}\SpecialCharTok{{-}}\DecValTok{12}\NormalTok{, }\DecValTok{1}\NormalTok{), }\AttributeTok{display\_xlim =} \FunctionTok{c}\NormalTok{(}\DecValTok{1980}\NormalTok{, }\DecValTok{2010}\NormalTok{))}
   \FunctionTok{points}\NormalTok{(data}\SpecialCharTok{$}\NormalTok{years[idxs\_t], data}\SpecialCharTok{$}\NormalTok{log\_rw\_obs[idxs\_t], }\AttributeTok{pch=}\DecValTok{16}\NormalTok{, }\AttributeTok{cex=}\FloatTok{0.7}\NormalTok{, }\AttributeTok{col=}\StringTok{"white"}\NormalTok{)}
  \FunctionTok{points}\NormalTok{(data}\SpecialCharTok{$}\NormalTok{years[idxs\_t], data}\SpecialCharTok{$}\NormalTok{log\_rw\_obs[idxs\_t], }\AttributeTok{pch=}\DecValTok{16}\NormalTok{, }\AttributeTok{cex=}\FloatTok{0.4}\NormalTok{, }\AttributeTok{col=}\StringTok{"black"}\NormalTok{)}
\NormalTok{\}}
\end{Highlighting}
\end{Shaded}

\includegraphics{marginal_shocks_files/figure-pdf/unnamed-chunk-14-1.pdf}

\linespread{2}

\newpage

However, it seems that the shocks are bit too large\ldots{}

\linespread{0.5}

\begin{Shaded}
\begin{Highlighting}[]
\FunctionTok{par}\NormalTok{(}\AttributeTok{mfrow =} \FunctionTok{c}\NormalTok{(}\DecValTok{2}\NormalTok{,}\DecValTok{2}\NormalTok{), }\AttributeTok{mar =} \FunctionTok{c}\NormalTok{(}\DecValTok{2}\NormalTok{, }\DecValTok{5}\NormalTok{, }\DecValTok{3}\NormalTok{, }\DecValTok{1}\NormalTok{))}
\ControlFlowTok{for}\NormalTok{(t }\ControlFlowTok{in}\NormalTok{ random\_trees)\{}
\NormalTok{  idxs\_t }\OtherTok{\textless{}{-}}\NormalTok{ data}\SpecialCharTok{$}\NormalTok{tree\_start\_idxs[t]}\SpecialCharTok{:}\NormalTok{data}\SpecialCharTok{$}\NormalTok{tree\_end\_idxs[t]}
\NormalTok{  names }\OtherTok{\textless{}{-}} \FunctionTok{sapply}\NormalTok{(idxs\_t,}
                  \ControlFlowTok{function}\NormalTok{(x) }\FunctionTok{paste0}\NormalTok{(}\StringTok{\textquotesingle{}log\_rw\_pred[\textquotesingle{}}\NormalTok{, x, }\StringTok{\textquotesingle{}]\textquotesingle{}}\NormalTok{))}
\NormalTok{  util}\SpecialCharTok{$}\FunctionTok{plot\_conn\_pushforward\_quantiles}\NormalTok{(samples, names, data}\SpecialCharTok{$}\NormalTok{years[idxs\_t], }
                                       \AttributeTok{baseline\_values =}\NormalTok{ data}\SpecialCharTok{$}\NormalTok{log\_rw\_obs[idxs\_t],}
                                       \AttributeTok{xlab=}\StringTok{""}\NormalTok{, }\AttributeTok{ylab=}\FunctionTok{paste0}\NormalTok{(}\StringTok{"Residuals {-} tree "}\NormalTok{,t), }
                                       \AttributeTok{display\_ylim=}\FunctionTok{c}\NormalTok{(}\SpecialCharTok{{-}}\DecValTok{3}\NormalTok{, }\DecValTok{1}\NormalTok{), }\AttributeTok{display\_xlim =} \FunctionTok{c}\NormalTok{(}\DecValTok{1980}\NormalTok{, }\DecValTok{2010}\NormalTok{),}
                                       \AttributeTok{residual =} \ConstantTok{TRUE}\NormalTok{)}
\NormalTok{\}}
\end{Highlighting}
\end{Shaded}

\includegraphics{marginal_shocks_files/figure-pdf/unnamed-chunk-15-1.pdf}

\linespread{2}

And, let's not forget the stand-level GP, which is still wiggly:

\linespread{0.5}

\begin{Shaded}
\begin{Highlighting}[]
\NormalTok{s }\OtherTok{\textless{}{-}} \DecValTok{1}
\NormalTok{names }\OtherTok{\textless{}{-}} \FunctionTok{sapply}\NormalTok{(}\DecValTok{1}\SpecialCharTok{:}\NormalTok{data}\SpecialCharTok{$}\NormalTok{N\_all\_years,}
                  \ControlFlowTok{function}\NormalTok{(y) }\FunctionTok{paste0}\NormalTok{(}\StringTok{\textquotesingle{}f\_sh[\textquotesingle{}}\NormalTok{, s, }\StringTok{\textquotesingle{},\textquotesingle{}}\NormalTok{, y, }\StringTok{\textquotesingle{}]\textquotesingle{}}\NormalTok{))}
\NormalTok{util}\SpecialCharTok{$}\FunctionTok{plot\_realizations}\NormalTok{(samples, names, }\AttributeTok{plot\_xs =}\NormalTok{ data}\SpecialCharTok{$}\NormalTok{all\_years[}\DecValTok{1}\SpecialCharTok{:}\NormalTok{data}\SpecialCharTok{$}\NormalTok{N\_all\_years], }\AttributeTok{N\_plots =} \DecValTok{50}\NormalTok{,}
                       \AttributeTok{main =} \FunctionTok{paste0}\NormalTok{(}\StringTok{\textquotesingle{}Stand \textquotesingle{}}\NormalTok{, data}\SpecialCharTok{$}\NormalTok{uniq\_stand\_ids[s]), }\AttributeTok{ylab =} \StringTok{\textquotesingle{}f\_sh\textquotesingle{}}\NormalTok{,}
                       \AttributeTok{display\_xlim =}\NormalTok{ data}\SpecialCharTok{$}\NormalTok{all\_years[}\FunctionTok{c}\NormalTok{(}\DecValTok{1}\NormalTok{,data}\SpecialCharTok{$}\NormalTok{N\_all\_years)], }\AttributeTok{display\_ylim =} \FunctionTok{c}\NormalTok{(}\SpecialCharTok{{-}}\DecValTok{2}\NormalTok{,}\DecValTok{2}\NormalTok{))}
\end{Highlighting}
\end{Shaded}

\includegraphics{marginal_shocks_files/figure-pdf/unnamed-chunk-16-1.pdf}

\linespread{2}




\end{document}
